\section{Special classes of geodesic currents}
\label{sec:special-classes}

In this section we characterize special classes of geodesic currents as curve functionals, such as filling currents, hyperbolic metrics, measured laminations, integral multi-curves, or currents dual to embedded graphs.

\subsection{Filling currents}
\label{subsec:filling}
Using a result of Burger, Iozzi, Parreau, and Pozzetti, we can give a
characterization of filling geodesic currents.

\begin{corollary}
  Let $S$ a closed surface,
  $\mu$ a geodesic current and $f$ the functional on curves associated
  to $\mu$, i.e., so that $\mu_f=\mu$. The following subsets are in
  bijection.

\begin{enumerate}
    \item The set of filling currents, i.e., $ \Curr_{\mbox{fill}}(S)$
    \item The set of functionals $f \colon \mathcal{C}^+(S) \to
      \mathbb{R}$ which satisfy symmetry, oriented smoothing,
      additivity, and stability, and so that their continuous extension to geodesic currents is positive.
    \item The set of curve functionals $f \colon \mathcal{C}(S) \to
      \mathbb{R}$ for which their systole \[\sys(f)\coloneqq \inf \{
        f(C) : C \in \mathcal{C}(S),\ C \mbox{ non-trivial} \}\] is
      positive and satisfy symmetry, oriented smoothing, additivity,
      and stability.
\end{enumerate}
\end{corollary}
\begin{proof}
These equivalences follow from Theorem~\ref{thm:intersect} and \cite[Theorem~1.3]{BIPP21:Systoles}.
\end{proof}

\begin{remark}
Since the space of geodesic currents is locally compact (see Proposition~\ref{prop:local_compactness}), this shows $\AConvex$ is locally compact in the product topology, which is not obvious.
See also Section~\ref{subsubsec:convexity_charts} for a related discussion.
\end{remark}

\subsection{Multi-curves}

We have the following characterization of curves that are one-half of an integral multi-curve.

\begin{corollary}
Let $f \in \AConvex(S)$ such that $f(C) \in \mathbb{Z}$ for every $C \in \mathcal{C}(S)$. Then there exists an integral multi-curve $M$ so that
\[
i(1/2 M,\cdot)=f(\cdot).
\]

\label{cor:int_multicurve}
\end{corollary}
\begin{proof}
Since $f \in \AConvex(S)$, by Theorem~\ref{thm:intersect} there exists a geodesic current $\mu$ so that
\[
i(\mu, C) = f(C)
\]
for every closed curve $C$.
It suffices to show $\mu$ is an integral multi-curve. By the same argument in the proof of~\cite[Th\'eor\`eme~1]{Otal90:SpectreMarqueNegative}, there exists a family of rectangles of geodesics $\mathcal{R}$ with dense endpoints in $\partial_{\infty} S$ whose measure can be approximated by a sequence of one-half of the difference of sums of intersection numbers. That is, for every $R \in \mathcal{R}$, there exist sequences $A_n, B_n, C_n, D_n \in \mathcal{C}(S)$ so that
\[
\mu(R) = \frac{1}{2}\lim_n i(\mu, B_n) + i(\mu, C_n) - i(\mu, A_n) - i(\mu, D_n).
\]
Since the expression in the limit is one-half of an integer number for every $n$, for $n$ large enough the sequence must be constant (by discreteness of $\mathbb{Z}$).
Hence, $\mu(R)$ is one-half of an integer for every $R \in \mathcal{R}$.
By~\cite[Proposition~4.12]{BIPP24:PositiveCrossratios}, it follows that $2\mu$ must be an integral multi-curve.
\end{proof}

\begin{remark}
  The factor of $1/2$ is necessary, since there are examples of
  word-lengths (whose marked length spectrum is integral) that are
  dual to one-half of an integral
  multi-curve~\cite[Theorem~1.1]{Erlandsson:WordLength}. For another
  concrete example, let $\mu$ be the geodesic current supported on a
  separating simple closed curve. Then $\mu/2$ has an integral
  length spectrum but is not an integral multi-curve.
\end{remark}

\DPTtodo[inline]{Characterize which half-integral
  multi-curves have integral intersection number using $\bbZ/2$
  homology. Also characterize intersection number with integral
  curves. DMG on it}

It is easy to see that if $M = \sum_{i=1}^n \lambda_i C_i$, where $\lambda_i \geq 0$ are non-negative real numbers, then for every closed curve $C$, there exist integer numbers $n_i$ so that
\begin{equation}
i(M, C) = \sum_{i=1}^n \lambda_i n_i.
\label{eq:weighted_multicurve}
\end{equation}

This property does not seem to characterize weighted multi-curves (see~\cite{Agol25:Spectrum}).

\subsection{Embedded graphs}
\label{subsec:embedded_graphs}
Related to multi-curves, we also obtain geodesic currents dual to stable lengths associated to embedded graphs in surfaces.

 Let $\iota \colon \mathcal{G} \to S$ be an \emph{embedded graph}, i.e., an embedding of a finite graph in $S$ that is filling, in the sense that the complementary regions are disks, or equivalently $i_*$ is surjective on $\pi_1$.
 Endow $\mathcal{G}$ with a length metric $g$. Then any closed multi-curve $C$ on $S$ can be homotoped so that it factors through $\mathcal{G}$, in many different ways. Let $\ell_{\mathcal{G}}(C)$ be the length of the smallest
 multi-curve $D$ on $\mathcal{G} $ so that $\iota(D)$ is homotopic to $C$. It is easy to see that this length is realized and is positive. 

\begin{theorem}
Let $\iota \colon \mathcal{G} \to S$ be an embedded graph. Then $\ell_{\mathcal{G}}$ is dual to a geodesic current $\mu_{\mathcal{G} }$.
\label{thm:dual_embeddedgraph}
\end{theorem}
\begin{proof}
In~\cite[4.4]{MGT21:Smoothings} we show that such lengths are in $\AConvex(S)$, hence by Theorem~\ref{thm:intersect} they have a dual geodesic current. 
\end{proof}

If furthermore $\mathcal{G}$ is an embedded graph with the edge-metric on $\pi_1(S)$ (i.e., each edge has weight 1), $\mu_{\mathcal{G}}$ must be $1/2$ of an integral multi-curve. This generalizes the main result in~\cite{Erlandsson:WordLength}, where Erlandsson finds that word-metric arising from \emph{simple generating set} are dual to 1/2 of an integral multi-curve dual. Such generating sets correspond to choosing $\mathcal{G}$ to be a rose graph with only one vertex and edges of length 1.

\begin{corollary}
Let $\mathcal{G}$ be an embedded graph with the edge-metric on $\pi_1(S)$ (i.e., each edge has weight 1). Then $\ell_{\mathcal{G}}$ is dual to a geodesic current $\mu_{\mathcal{G}}$ which is $1/2$ of an integral multi-curve.
\label{cor:generalization_wordlength}
\end{corollary}
\begin{proof}
Since $\mathcal{G}$ is equipped with the edge-metric, $\ell_{\mathcal{G}}(C) \in \mathbb{Z}$.
Hence, Corollary~\ref{cor:int_multicurve} implies the result.
\end{proof}

We note that our result says nothing explicit about the dual multi-curve, whereas Erlandsson's result explicitly constructs it in the specific case of simple generating sets. It would be interesting to describe these multi-curves in general.

\subsection{Measured laminations}

We will now characterize measured laminations as curve functionals in $\AConvex(S)$.

We will need an observation about the measure of sets intersecting pairs of geodesic segments, proven in \cite[Lemma~6.2]{dRMG22:Duals} in the setting of unoriented geodesic currents. We will prove an analogous oriented version.

We introduce the following notation. Given two geodesics segments $t$ and $t'$, let $G(t , t')$ be the set of geodesics intersecting from segment $t$ to segment $t'$ transversely. We will refer to sets of this type as \emph{double transversals}.

Let $x,y,z,w \in \wt{X}$ appear counter-clockwise as vertices of an
embedded convex geodesic quadrilateral on $\wt{X}$.

\begin{lemma}
Given the setup described above, we have 
\begin{align*}
\mu(G([x,w), [y,z))) + \mu(G((x,w], (y,z])) &=
\mu(G[x,z)) + \mu(G[w,y))
- \mu(G([w,z)) - \mu(G([x,y))\\
\mu(G((x,w], (y,z])) + \mu(G([x,w), [y,z))) &=
\mu(G(x,z]) + \mu(G(w,y])
- \mu(G((w,z]) - \mu(G((x,y])
\end{align*}
\label{lem:otalclaim}
\end{lemma}

We also need the following hyperbolic geometry lemma.

\begin{proposition}
Consider two elements $a, b \in \pi_1(\Sigma)$ with hyperbolic axes $\ora{a}$ and $\ora{b}$ that intersect at a point $O$, and translation lengths $\ell(a),\ell(b)$.
We define the following four points:

\begin{itemize} \item $P$ is the point on $\ora{a}$ obtained by moving forward from $O$ by distance $\ell(a)/2$;
\item $Q$ is the point on $\ora{b}$ obtained by moving backward from $O$ by distance $\ell(b)/2$;
\item $Q'$ is the point on $\ora{b}$ obtained by moving forward from $O$ by distance $\ell(b)/2$;
\item $P'$ is the point on $\ora{a}$ obtained by moving backward from $O$ by distance $\ell(a)/2$.
\end{itemize}
We have the following (see Figure~\ref{fig:smoothing_1}):
\begin{enumerate}
\item The segment $(Q,P]$ is on the line $\ora{ab}$ and has length $\ell(ba)/2$.
\item The segment $(P',Q']$ is on the line $\ora{ba}$ and has length $\ell(ba)/2$.
\item The segment $(Q',P]$ is on the line $\ora{ab^{-1}}$ and has length $\ell(ab^{-1})/2$.
\item The segment $(P',Q]$ is on the line $\ora{b^{-1}a}$ and has length $\ell(ab^{-1})/2$.
\end{enumerate}
Furthermore, $G((Q,P])\cup a\cdot
G((P',Q'])$ is a Borel fundamental domain for the
action of $\langle ab \rangle$ on $G(\ora{ab})$.
Similarly, $G((P,Q']) \cup b \cdot G((Q,P'])$ is a Borel fundamental domain for
 the action of $\langle bA\rangle$ on $G(\ora{bA})$.
\label{prop:keygeometry}
\end{proposition}

  \begin{figure}
\centering{
\svgfignormal{smoothing_1}}
\caption{Hyperbolic geometry of the smoothing between two distinct curves.}
\label{fig:smoothing_1}
\end{figure}

Figure~\ref{fig:smoothing_1} shows a \emph{hyperbolic parallelogram},
described more in Appendix~\ref{sec:hyp-par}.

\begin{proof}
  We prove the first case.
  For $z$ any of the five points in the statement, let $r_z$ be
  the hyperbolic rotation by angle $\pi$ fixing~$z$. Then $a$, as
  acting on the hyperbolic plane, is $r_P r_O$, while $b = r_O r_Q$.
  (Recall that a composition $r_zr_w$ of two rotations by $\pi$ is a
  translation on the line $\overleftrightarrow{zw}$ by twice the
  distance between $z$ and~$w$.)
  It follows that $ab = r_P (r_O)^2 r_Q =
  r_Pr_Q$, from which the statement about $(Q,P]$
  follows. The other cases are similar.

  \begin{comment}
  Since $\ora{a}$ and $\ora{b}$ intersect at $O$, so do the curves $[a]$ and $[b]$ at some point $o \in \Sigma$.
Perform a smoothing operation obtaining a new curve $\alpha$ of length $\ell(a) + \ell(b)$.
A lift $\alpha$  of this curve in $\wt{X}$ going through $O$ is a piecewise geodesic consisting of a segment on $\ora{b}$ of length $\ell(b)$, then a segment of $\ora{b^{-1}ab}$ of length $\ell(a)$, etc.
Let $\alpha_*$ be the geodesic representative in the homotopy class of $\alpha$.
Now, connecting the midpoints of the geodesic segments of $\alpha$ we obtain another piecewise geodesic $\alpha'$ which is, in fact, a geodesic segment.
Indeed, we note that $\alpha$ is invariant under angle $\pi$ rotation around any of the intersection points between $\alpha$ and $\alpha'$.
Since $\alpha$ and $\alpha'$ have the same stabilizer, they are lifts of the same geodesic $\alpha$.
Thus, the lift of $\alpha_*$ is $\alpha'$, which is determined by two of the midpoints defining $\alpha$, such as $P'$ and $Q'$.
As for the translation lengths, we note that $b=r_O r_{Q'}$ where $r_O$ denotes the reflection along the geodesic perpendicular to $\ora{b}$ at $O$,
and $r_{Q'}$ is the corresponding reflection at $Q'$.
Similarly, we can write $a=r_{P'}r_O$, and thus $ab=(r_{P'}r_O)(r_O r_{Q'})=r_{P'}r_{Q'}$, which corresponds to a translation along $\alpha'$ of translation length
twice the distance between $P'$ and $Q'$.
The other cases follow similarly.
\end{comment}

For the assertion on fundamental domains, we note that the segment
$(P',Q']$ on the axis of $\ora{ba}$, is sent by $a$ to a segment
$(P,R]$, where $R=a(Q')$ appears forward of $P$ in according to the
orientation of $\ora{ab}$, and the segment $(Q,R]$ is a fundamental
domain for the action of $ab$ on its axis $\ora{ab}$. Thus $\{
\ora{ab} \} \times G((Q,R])= \{ \ora{ab} \} \times  (G((Q,P]) \cup a
\cdot G((P',Q']))$ is a fundamental domain for the action of $ab$ on
$\DG(\Sigma)$, as desired.
The other claim is similar.
\end{proof}

The above lemma can be seen as a refinement of~\cite[Lemma~5.6]{NCS23:Property}.
Using it, we characterize measured laminations among curve
functionals.

\DPTtodo[inline]{I believe we \emph{do} need to be careful about
  endpoints, since we can have non-zero atoms there, and we use
  max-smoothing for each $n$. This also applies to $B_{g,h}$, since
  the $\overline{R_n}$ intersect the boundary of $\overline{B_{g,h}}$.
  See eg the note about the backwards implication}
\DPTtodo[inline]{There's also a lot of poorly defined notation, eg the
  $\perp$ variations are introduced in a scattershot way. Waiting to
  fix that until we do the futzing with endpoints}
\begin{theorem}
  A curve function $f$ is dual to a measured lamination if and only if
  $f$ satisfies the hypotheses of Theorem~\ref{thm:intersect} and the
  \emph{disconnected max-smoothing property}:
    For any pair of crossing elements $g,h \in \pi_1(S)$, we have
    \begin{equation}\label{eq:disconn-max-smoothing}
      f\left([g][h]\right) = \max \left\{ f\left([gh]\right),f\left([gh^{-1}]\right) \right\}.
    \end{equation}
    \label{thm:char_laminations}
\end{theorem}
\begin{proof}
Suppose first $f$ satisfies Eq.~\eqref{eq:disconn-max-smoothing}.
By Theorem~\ref{thm:intersect}, there exists a geodesic current $\mu_f$ so that $f(C)=i(\mu_f,C)$ for all closed curve $C$.
It remains to show that $\mu_f$ is a measured lamination.
Given elements $g, h \in \pi_1(\Sigma)$ with intersecting axes,
consider the hyperbolic parallelogram~$H$ at the center of
Figure~\ref{fig:smoothing_1}.
Consider the set $R_{a,b}^+ \coloneqq G([Q,P) \times [P', Q'))$ and  $R_{a,b}^- \coloneqq G((Q,P] \times (P', Q'])$. Barrying endpoints, these two sets roughly consist of geodesics intersecting both sides of $H$ contained in $\ora{hg}$ and $\ora{gh}$.
We will also define, analogously, the sets
$R_1^{\perp}$

the set of geodesics intersecting the sides of $H$ on the axes $\ora{gh}$ and $\ora{hg}$.
\DPTtodo{Ambiguous: there are two natural boxes here.}
Also let $B_{g,h}$ be the box with vertices the endpoints of $g$ and $h$.

Taking powers $g^n$ and $h^n$, and we obtain corresponding hyperbolic
parallelograms $H_n$ for $g^n,h^n$ and analogous sets $R_n \coloneqq
R_{g^n,h^n}$.
Observe that $R_n \subset R_{n+1}$, and $\bigcup_{n \geq 1} R_n = B_{g,h}$.
Hence we have
\[
\lim_{n\to\infty} \mu(R_n) \to \mu(B_{g,h}).
\]
If we let $a_n=g^n$, $b_n=h^n$, $c_n=a_n b_n^{-1}$ and $d_n=a_n b_n$,
we get, by Lemma~\ref{lem:otalclaim} and Proposition~\ref{prop:keygeometry},
\DPTtodo{I think $C_n$ and $D_n$ might be switched in these equations.
  Also it seems this is an independent lemma, also used in
  Theorem~\ref{thm:char_liouville}}
\begin{equation}\label{eq:Rn-meas}
\mu(R_n)=A_n + B_n-C_n,
\end{equation}
where $A_n=f(a_n)/2$, $B_n=f(b_n)/2$, $C_n=f(a_nb_n^{-1})/2$.
If we further let $R_n^{\perp}$ be the set of geodesics intersecting the sides of $H_n$ contained in $\ora{a_nb_n^{-1}}$ and $\ora{b_n^{-1}a_n}$, we have
\begin{equation}\label{eq:Rn-perp-meas}
\mu(R_n^\perp)=A_n + B_n -D_n,
\end{equation}
where $D_n=f(a_nb_n)/2$.
We also have $R_n^{\perp} \subset R_{n+1}^{\perp}$ and $\bigcup_{n \geq 1} R_n^{\perp} = B_{g,h}^{\perp}$, where $B_{g,h}^{\perp}$ denotes the box complementary to $B_{g,h}$.
Thus, as above, we have $\lim_n \mu(R_n^{\perp})=\mu(B_{g,h}^{\perp})$.
By max-smoothing, for every $n$, we have $\mu(R_n^\perp)=0$ or $\mu(R_n)=0$.
Since each sequence of sets $(R_n)$ and $(R_n^{\perp})$ is nested, if $\mu(R_1) \neq 0$, then $\mu(R_n) \neq 0$ for all $n \geq 1$ and thus $\mu(R_n^{\perp})=0$ for all $n \geq 1$.
Hence, $\mu(B_{g,h})=0$ or $\mu(B_{g,h}^{\perp})=0$.
By density of endpoints of closed geodesics, these boxes have vertices on a dense set of points of $\partial_{\infty} S$.
By density of $B_{g,h}$ in the space of boxes of geodesics as $g,h$ runs over intersecting pairs in $\pi_1(\Sigma)$, given any box $B$, we can find a sequence of intersecting pairs $g_n, h_n \in \pi_1(S)$ so that if we let $B_n \coloneqq B_{g_n, h_n}$, we have $B_n \subset B_{n+1}$, $B_{n+1}^{\perp} \subset B_{n}^{\perp}$, $\bigcup_n B_n = B$, and $\bigcap_n B_n^{\perp} = B$.
Then, by lower and upper continuity of measures, respectively, we get $\lim_n \mu(B_n) =\mu(B)$ and $\lim_n \mu(B_n^{\perp}) = \mu(B)$.

Now we prove the other implication, by proving the contra-positive.
If disconnected max-smoothing is \emph{not} satisfied for the elements
$g,h$, then by Equations~\eqref{eq:Rn-meas}
and~\eqref{eq:Rn-perp-meas}, both $\mu(R_{g,h})>0$ and
$\mu(R_{g,h}^{\perp})>0$.
\DPTtodo{I believe appropriately interpreted neither $R_{g,h}$ nor
  $R_{g,h}^\perp$ contain the diagonals $\ora{g}$ and $\ora{h}$. We
  use that here}
Since for
every $\gamma \in R_{g,h}$ and every $\gamma' \in R_{g,h}^\perp$, the geodesics
that $\gamma$ and $\gamma'$ cross, this implies that $i(\mu, \mu) >0$,
and hence $\mu$ is not a measured lamination.
\end{proof}

The above result shows that if one knows $f$ is dual to a geodesic
current, then it is only necessary to check disconnected max-smoothing
to conclude $f$ is a measured lamination, i.e., one does not have to
check essential self-crossings. It is however true that intersection
number with a measured lamination also satisfies connected
max-smoothing; this is Theorem~\ref{mainthm:apps}\ref{item:lam-char}.

\begin{theorem}
A curve functional $f$ is dual to a measured lamination if and only if it satisfies additivity, stability and the \emph{max-smoothing property}:
    For any essential crossing of a multi-curve $C$, we have
    \begin{equation}\label{eq:MLsmoothing}
      f\left(\mfig{curves-0}\right) = \max \left\{ f\left(\mfig{curves-1}\right),f\left(\mfig{curves-2}\right) \right\}.
    \end{equation}
    \label{prop:char_laminations_plus}
\end{theorem}
\begin{proof}
Suppose $f$ is a curve functional satisfying max-smoothing, additivity and stability. Max-smoothing clearly implies smoothing, as well as disconnected max-smoothing. Thus, by Theorem~\ref{thm:char_laminations}, $f$ is dual to a measured lamination. On the other hand, $f$ is dual to a measured lamination,  Theorem~\ref{thm:char_laminations} implies disconnected max-smoothing. Connected max-smoothing follows from tracing through the proof of Proposition~\ref{prop:cont_conn_disc}, using the fact that $i(\mu, \cdot)$ is continuous, additive and stable, and noting that if $f$ satisfies disconnected smoothing, all the inequalities become equalities in the proof of disconnected smoothing implies connected smoothing.
\end{proof}

The fact that measured laminations satisfy max-smoothing was proven by
the second author in~\cite{Thu00:GeometricIntersection} using
different techniques. A related characterization of measured
laminations in terms of valuations in the character variety has been
proven recently by Marché and Simon~\cite[Theorem~3.1]{MS21:Valuations}.

\subsection{Liouville currents}
\label{subsec:liouville}
Bonahon gives an algebraic characterization for a geodesic current to be a Liouville current for a hyperbolic structure $X$.

\begin{theorem} [{\cite[Theorem~13]{Bonahon88:GeodesicCurrent}}]
   A geodesic current $\mu \in \GC(S)$ is $\mu=\mathcal{L}_{\Sigma}$ for some $\Sigma$ if and only if, for any four points $x,y,z,w \in \partial_{\infty} S$, we have
   \[
      e^{-\mu([x,y]\times[z,w])} + e^{-\mu([y,z]\times [w,x])}=1.
   \]
   \label{thm:bonahonliouville}
\end{theorem}
\DPTtodo[inline]{Do we need to worry about atoms in above statement?}

We use this to prove the following characterization (Theorem~\ref{mainthm:apps}\ref{item:hyp-char}).

\begin{theorem}
  A functional $f$ is dual to a hyperbolic Liouville current if and only
  if it satisfies additivity, stability, and,
    for any pair of crossing elements $g, h \in \pi_1(S)$ we have the \emph{$\lambda$-smoothing relation}  \[
        \lambda \left(g\right)\lambda \left(h\right) = \lambda \left(gh\right) + \lambda \left(gh^{-1}\right),
  \]
  where $\lambda( \cdot ) \coloneqq 2\cosh(f(\cdot)/2)$.
    \label{thm:char_liouville}
  \end{theorem}
\begin{proof}
The proof is similar to that of Theorem~\ref{thm:char_laminations}.

\DPTtodo[inline]{Moved the easy direction first}
If $\mu$ is hyperbolic Liouville current then the disconnected $\lambda$-smoothing property follows from the relation between the hyperbolic lengths of the sides of a hyperbolic parallelogram (see Lemma~\ref{lem:hyp_parallelogram}).


\DPTtodo[inline]{This seems to be missing the short argument that
  the eqn above implies smoothing. (In
  Theorem~\ref{thm:char_liouville_strong} a more precise version is
  invoked.) Also, we invoke disconnected
  $\Rightarrow$ connected smoothing
  here, right?}
For the other direction, by the assumptions on $f$, there exists a
geodesic current $\mu_f$ so that $f(C)=i(\mu_f,C)$ for all closed
curve $C$.
It remains to show that $\mu_f$ is a Liouville current.
Given elements $g, h \in \pi_1(\Sigma)$ with $\ora{g}, \ora{h}$ crossing, we consider powers $g^n$ and $h^n$, and obtaining a corresponding hyperbolic parallelogram $H_n$ for $g^n,h^n$, we obtain the set $R_n \coloneqq R_{g^n,h^n}$ defined as in the proof of Theorem~\ref{thm:char_laminations}.
By the same considerations as in that proof, letting $a_n=g^n$, $b_n=h^n$, $c_n=a_n b_n^{-1}$ and $d_n=a_n b_n$,
we get
\[\mu(R_n)=A_n + B_n-C_n\]
and
\[
\mu(R_n^\perp)=A_n + B_n -D_n,
\]
where $A_n=f(a_n)/2$, $B_n=f(b_n)/2$, $C_n=f(a_nb_n^{-1})/2$ and $D_n=f(a_nb_n)/2$.
Now, we compute
\[
e^{-\mu(R_n)} + e^{-\mu(R_n^{\perp})} = e^{-(A_n+B_n)+C_n} + e^{-(A_n+B_n)+D_n} = e^{-(A_n+B_n)} (e^{C_n} + e^{D_n}).
\]
On the other hand, when $n$ goes to infinity, we have $e^{A_n} \sim \lambda(a_n)$, $e^{B_n} \sim \lambda(b_n)$, $e^{C_n} \sim \lambda(c_n)$ and $e^{D_n} \sim \lambda(d_n)$ (where $\sim$ means that the ratio of the left hand side to the right hand side tends to 1 as $n$ tends to infinity). Since the disconnected $\lambda$-smoothing relation for $a_n, b_n$ is given by
\[
\lambda(a_n)\lambda(b_n) = \lambda(c_n) + \lambda(d_n),
\]
and a straightforward computation shows $\lambda(c_n) + \lambda(d_n) \sim e^{C_n}+e^{D_n}$, and we get
\[
e^{A_n}e^{B_n} \sim e^{C_n}+e^{D_n}.
\]
Therefore, we have
\[
\lim_{n\to\infty} e^{-\mu(R_n)} + e^{-\mu(R_n^{\perp})}=1.
\]
We also have $R_n^{\perp} \subset R_{n+1}^{\perp}$ and $\bigcup_{n \geq 1} R_n^{\perp} = B_{g,h}^{\perp}$, where $B_{g,h}^{\perp}$ denotes the box complementary to $B_{g,h}$, and $B_{g,h}$ is the open box determined by the endpoints of $g$ and $h$.
Thus, we have $\lim_n \mu(R_n)=\mu(B_{g,h})$ and $\lim_n \mu(R_n^{\perp})=\mu(B_{g,h}^{\perp})$.
Therefore, we have the relation
\[
e^{-\mu(B_{g,h})} + e^{-\mu(B_{g,h}^{\perp})} =1
\]
for any intersecting pair $g,h \in \pi_1(\Sigma)$.
As in the previous proof, by density of $B_{g,h}$ in the space of boxes of geodesics as $g,h$ runs over intersecting pairs in $\pi_1(\Sigma)$, given any box $B$, we can find a sequence of intersecting pairs $g_n, h_n \in \pi_1(S)$ so that if we let $B_n \coloneqq B_{g_n, h_n}$, we have $\lim_n \mu(B_n) =\mu(B)$ and $\lim_n \mu(B_n^{\perp}) = \mu(B)$, and hence
\[
e^{-\mu(B)} + e^{-\mu(B^{\perp})}=1.
\]
Therefore, Theorem~\ref{thm:bonahonliouville} shows $\mu$ is a Liouville current.
\end{proof}

Bonahon's Theorem~\ref{thm:bonahonliouville} was strengthened by
Hamenst\"adt

\begin{theorem}[{\cite[Corollary~2.12]{Ham99:Cocycles}}]
   A geodesic current $\mu \in \GC(S)$ is $\mu=\mathcal{L}_{\Sigma}$ for some $\Sigma$ if and only if, there exist a number $\rho>1$ so that for any four ordered points $x,y,z,w \in \partial_{\infty} S$, we have
    \[
      \frac{1}{\rho} \cdot  e^{-\mu([x,y]\times[z,w])} \leq 1- e^{-\mu([y,z]\times [w,x])} \leq \rho \cdot e^{-\mu([x,y]\times[z,w])} 
   \]
   \label{thm:ham}
\end{theorem}

Using this, we can prove the following.

\begin{theorem}
A geodesic current $\mu \in \GC(S)$ is a hyperbolic Liouville current if and only if the functional $f$ corresponding to $\mu$ satisfies smoothing, additivity, stability, and the following inequalities.
    There exists a constant $C \geq 1$ so that, for any pair of \emph{crossing} elements $g,h \in \pi_1(S)$, we have
    \begin{gather}
    1/C \leq \exp \Bigl( f\left([g][h]\right)/2-f\left([gh]\right)/2 \Bigr) - \exp \Bigl( f\left([gh^{-1}]\right)/2-f\left([gh]\right)/2 \Bigr) \leq C\\
     1/C \leq  \exp \Bigl( f\left([g][h]\right)/2-f\left([gh^{-1}]\right)/2 \Bigr) - \exp \Bigl( f\left([gh]\right)/2-f\left([gh^{-1}]\right)/2 \Bigr) \leq C.
    \end{gather}
    \label{thm:char_liouville_strong}
  \end{theorem}

\begin{proof}
The assumptions can be written as 
 \begin{gather}
     1/C \leq  \frac{1- \exp \left[f\left([gh]\right)/2 - f\left([gh^{-1}]\right)/2 \right]}{\exp \left[ -f\left([g][h]\right)/2 + f\left([gh]\right)/2 \right]} \leq C\\
    1/C \leq \frac{1- \exp \left[f\left([gh^{-1}]\right)/2 - f\left([gh]\right)/2\right]}{\exp \left[-f\left([g][h]\right)/2 + f\left([gh^{-1}]\right)/2\right]} \leq C.
\end{gather}
\DPTtodo{I feel like there are some steps missing. Does this use the
  $A_n,B_n,\dots$?}
    Now, as in the proof of Theorem~\ref{thm:char_liouville}, this gives us
     \begin{gather}
    1/C \leq \frac{1 - e^{-\mu(B_{g,h})}}{e^{-\mu(B_{g,h}^{\perp})}} \leq C\\
    1/C \leq \frac{1 - e^{-\mu(B_{g,h}^{\perp})}}{e^{-\mu(B_{g,h})}} \leq C
    \end{gather}
    and, again, by the same density argument in that proof, one gets
    the same inequalities for a general box of geodesics $B$. Hence,
    by Theorem~\ref{thm:ham}, we see that $\mu$ is a hyperbolic Liouville current.
On the other hand, if $\mu$ is hyperbolic Liouville current then the
disconnected $\lambda$-smoothing property is satisfied, and the
condition of the statement then follows by the same argument as in the
proof of Theorem~\ref{thm:char_liouville}
\DPTtodo{Did the $C=1$ argument get edited out?}
(which shows it is satisfied
with $C=1$).
\end{proof}

%%% Local Variables:
%%% mode: LaTeX
%%% TeX-master: "Intersections"
%%% End:
